\Subhead{1.11.2}{Возможные пути потери устойчивости регулярных колебаний и перехода к хаосу}
При изменении управляющего параметра --- например, накачки или глубины обратной связи --- устойчивое периодическое движение может потерять устойчивость. Но отдельная бифуркация ещё не означает появления хаоса: новым режимом может оказаться другой цикл или квазипериодическое движение. Важно проследить последовательность изменений установившегося движения.

\paragraph{Удвоения периода.}
Вместо одного повторяющегося оборота появляются два различающихся оборота: полное состояние повторяется через $2T$, затем возможны периоды $4T$, $8T$ и т.~д. На сечении Пуанкаре число точек последовательно удваивается; в спектре возникают субгармоники исходной частоты. Накопление таких бифуркаций способно привести к хаотическому режиму. Это один из возможных сценариев, а не обязательный путь любого генератора. Внутри хаотических областей могут снова появляться окна регулярных колебаний.

\paragraph{Перемежаемость.}
Почти периодические участки чередуются с нерегулярными всплесками. Система подолгу движется в области фазового пространства, где ранее существовал или был близок к устойчивости регулярный режим, затем быстро покидает её и возвращается. При изменении параметра меняются длительности регулярных участков и частота всплесков. Поэтому переход к хаосу может проявляться постепенно через редкие нарушения почти регулярного движения.

\paragraph{Через квазипериодическое движение.}
К исходной частоте добавляется независимая несоизмеримая частота; траектория уже не замыкается, а заполняет тор. Само это движение ещё регулярно: его спектр состоит из комбинационных линий. Дальнейшие бифуркации и разрушение тора могут привести к хаосу. Поэтому утверждение «колебания непериодические» недостаточно: необходимо различать квазипериодичность и хаотическую неустойчивость.
% Кузнецов Ю.И., §9.2, печатные с.112--117: три сценария.
% Теория универсальности Фейгенбаума и специальные модели статьи о КПР не включены.
